\documentclass[10pt]{article}
\usepackage[utf8]{inputenc}
\usepackage[T1]{fontenc}
\usepackage{amsmath}
\usepackage{amsfonts}
\usepackage{amssymb}
\usepackage[version=4]{mhchem}
\usepackage{stmaryrd}
\usepackage{bbold}

\begin{document}
а при $n=2 m$ равен

\[
\left(q^{2 m-1}-\varepsilon q^{m-1}\right)\left(q^{2 m-2}-1\right) q^{2 m-3} \ldots\left(q^{2}-1\right) q,
\]

где $\varepsilon=1$ при $(-1)^{m} \Delta \in \mathbb{F}_{q}^{2}$ и $\varepsilon=-1$ в противном случае. Указанные формулы вместе с приведенными общими фактами об О. г. при $v \geqslant 1$ позволяют вычислить также и порядки $\Omega_{n}$ и $P \Omega_{n}$, так как $v \geqslant 1$ при $n \geqslant 3$, а порядок $k^{*} / k^{*}$ равен 2. Группа $P \Omega_{n}, n \geqslant 5$, является одной из классических простых конечных групп (см. также III валле apynna).

Один из основных результатов об автоморфизмах О. г. состоит в следующем: если $n \geqslant 3$, то всякий автоморфизм $\varphi$ группы $O_{n}$ имеет вид $\varphi(u)=\chi(u) g u g^{-1}$, $u \in \hat{O}_{n}$, где $\chi-$ фиксированный гомоморфизм $O_{n}$ в ее центр, а $g$ - фиксированное биективное полулинейное отображение $V$ в себя, удовлетворяющее условию $Q(g(v))=r_{g} Q^{\sigma}(v)$ для всех $v \in V$, где $r_{g} \in k^{*}$, а $\sigma-$ связанныӥ с $g$ автоморфизм $k$. Если $v \geqslant 1$ и $n \geqslant 6$, то всякий автоморфизм $O_{n}^{+}$индуцирован автоморфизмом $O_{n}$ (см. [1], [3]).

Так же, как и другие классич. группы, О. г. допускает (при нек-рых предположениях) геометрич. характергвацию. А именно, пусть $Q$ - такая анизотропная форма, что $Q(v) \in k^{2}$ для любого $v \in V$. В әтом случаө $k-$ пифагорово упорядочиваемое поле. При фиксированном упорядочөнии поля $k n$-мерной ц е п в ю и н ц и д ө в тны $\mathbf{x}$ п о л у п р о с т $\mathrm{p}$ а н с т в в $V$ наз. любая последовательность $\left(H_{s}\right)_{1 \leqslant s \leqslant n}$, построенная по линейно независимой системе векторов $\left(h_{s}\right)_{1 \leqslant s \leqslant n}$, где $I_{s}-$ множество всех линейных комбинаций вида $\sum_{j=1}^{s} \lambda_{j} h_{j}, \lambda_{s} \geqslant 0$. Группа $O_{n}$ обладает с в о й с т в ом с в о б о д в о й по д в и ж н о с т и, т. е. для любых двух $n$-мерных цепей полупространств существует единственное преобразование из $O_{n}$, переводящее первую цепь во вторую. Это свойство характеризует О. г.: если $L-$ любое упорядоченное тело и $G$ - подгруппа в $G L_{n}(L), n \geqslant 3$, обладающая свойством свободной подвижности, то $L$ являөтся пифагоровым полем, а $G=O_{n}(L, f)$, где $f-$ такая анизотропная симметрич. билинейная форма, что $f(v, v) \in L^{2}$ для любого вектора $v$.

Пусть $\bar{k}$ - фиксированное алгебраич. замыкание поля $k$. Форма $f$ естественно продолжается до невырожденной симметрич. билинейной формы $\bar{f}$ на $V \otimes_{k} \bar{k}$, а О. г. $O_{n}(\bar{k}, f)$ является определенвой над $k$ линейной алгебраической әруппой с группой $k$-точек $O_{n}(k, f)$. Определяемые таким образом (для разных $f$ ) линейные алгебраич. группы изоморфны над $\bar{k}$ (но, вообще говоря, не над $k$ ); соответствующая линейная алгебраич. грушпа над $\bar{k}$ наз. о р т о г о наль н о й а лге и ч е с к о й г р у п п о й $O_{n}(\bar{k})$. Ее подгруппа $O_{n}^{+}(\bar{k}, \bar{f})$ также является линейной алгебраич. группой вад $\bar{k}$ ІІ наз. с о б с т в е н н о ор то г о н а л н о й, или c пеци альной ортогональной, ал ге бр а и ч е с к о й группой (обозначение: $S O_{n}(\bar{k})$ ); она является связной компонентой единицы группы $O_{n}(\bar{k})$. Группа $S O_{n}(\bar{k})$ - почти простая алгебраич. группа (т. е. не содержащая ненульмерных алгебраич. нормальных делителей) типа $B_{s}$ при $n=2 s+1, s \geqslant 1$, и типа $D_{s}$ при $n=2 s, s \geqslant 3$. Универсальной накрывающей группы $S O_{n}$ является спинорная группа.

Если $k=\mathbb{R}, \mathbb{C}$ или $p$-адическое поле, то $O_{n}(k, f)$ естественно снабжается структурой вещественной, комплексной или p-адической аналитической әруппы. $\Gamma$ рупша Ли $O_{n}(\mathbb{R}, f)$ определяется с точностью до ивоморфизма сигнатурой формы $f ;$ если эта сигнатура имеет вид $(p, q), p+q=n$, то $O_{n}(\mathbb{R}, f)$ обозначается через $O(p, q)$ и наз. п с е в д о о $\mathrm{p}$ т о г о н а л ь н о й $\mathbf{\Gamma} \mathbf{p} \mathbf{y}$ пц 0 й. Еө можно отождествить с группой Ли всех дейст- вительных $(n \times n)$-матрип $A$, удовлетворяющих условию

\[
A^{\top} I_{p, q} A=I_{p, q}, \quad \text { где } I_{p, q}=\left\|\begin{array}{cc}
1_{p} & 0 \\
0 & -1_{q}
\end{array}\right\|
\]

(через $1_{s}$ обозначена единичная $(s \times s)$-матрица); алгебра Ли этой группы отождествляется с алгеброй Ли всех действительных $(n \times n)$-матриі $X$, удовлетворяющих условию $X^{\top} I_{p, q}=-I_{p, q} X$.В частном случае $q=0$ группа $O(p, q)$ обозначается через $O(n)$ и няз. в е щ е с тв ен ной ортогональ но й алгебра Ли состоит из всех кососимметрических действительных $(n \times n)$-матриц. Группа Ли $O(p, q)$ имеөт четыре компоненты связности при $q \neq 0$ и две компоненты связности при $q=0$. Связной компонентой единицы является ее коммутант, к-рыӥ при $q=0$ совпадает с подгруппой $S O(n)$ в $O(n)$, состоящей из всех преобразований с определителем, равным 1. Группа $O(p, q)$ комтактна только при $q=0$. Инварианты $S O(n)$ как топологич. многообразия достаточно подробно изучены. Один из классич. результатов в әтом направлении - вычисление чисел Бетти многообразия $S O(n)$ : его полином Пуанкаре имеет вид

\[
\prod_{s=1}^{m}\left(1+t^{4 s-1}\right)
\]

при $n=2 m+1$ и вид

\[
\left(1+t^{2 m-1}\right) \prod_{s=1}^{m-1}\left(1+t^{4 s-1}\right)
\]

при $n=2 m$. Фундаментальная группа многообразия $S O(n)$ есть $\mathbb{Z}_{2}$. Вычисление высших гомотопич. групп $\pi_{l}(S O(n))$ имеет непосредственное отношение к классификации локально тривиальных главных $S O(n)$ расслоений над сферами. Важную роль в тонологической $K$-теории играет теорема периодичности, согласно к-рой при $N \gg n$ имеют место изоморфизмы

\[
\pi_{n+8}(O(N)) \simeq \pi_{n}(O(N)), \pi_{n}(O(N)) \simeq \mathbb{Z}_{2},
\]

если $n=0,1$;

если $n=3,7$, п

\[
\begin{aligned}
& \pi_{n}(O(N)) \simeq \mathbb{Z}, \\
& \pi_{n}(O(N))=0,
\end{aligned}
\]

если $n=2,4,5,6$. Изучение топологии группы $O(p, q)$ по существу сводится к предыдущему случаю, т. к. связная компонента единицы группы $O(p, q)$ диффеоморфина произведению $S O(p) \times S O(q)$ на евклидово пространство.

Лит.: [1] Д в е д о н н е Ж., Геометрия классических групп, пер. с франц., М., 1974; [2] А р т и н Э., Геометрическая алгебра, пер. с англ., М., 1989; [3] Автоморфизмы классических групп, пер. с англ. и франц., М., 1976; [4] В е й л ь Г., Класси$1947 ;$ [5] Ж е л о б е н к о Д. П., Компактные группы Ли и их представления, М., 1970; [6] Б у р б а к и Н., Алгебра. Модули, кольца, формы, пер. с франц., М., 1966; [7] O'M e a r a O. T., Introduction to quadratic forms, B.-Hdlb., 1963; [8] X $\mathrm{X}$. м о л л е р Д., Расслоенные пространства, пер. с англ., М., 1970.

ОРТОГОНАЛЬНАЯ МАТРИЦА - матрица над коммутативным кольцом $R$ с единицей 1 , для к-рой транспонированная матрица совпадает с обратной. Определитель О.м. равен \textbackslash pm 1 . Совокупность всех $О$. м. порядка $n$ над $R$ образует подгруппу полной линейной группы $\mathrm{GL}_{n}(R)$. Для любой действительной $\mathrm{O}$. м. $a$ существует такая действительная О. м. $c$, что

\[
c a c^{-1}=\operatorname{diag}\left[ \pm 1, \ldots, \pm 1, a_{1}, \ldots, a_{t}\right]
\]

где

\[
a_{j}=\left\|\begin{array}{rr}
\cos \varphi_{j} & \sin \varphi_{j} \\
-\sin \varphi_{j} & \cos \varphi_{j}
\end{array}\right\|
\]

Невырожденная комплексная матрица $a$ тогда і только тогда подобна комплексной $\mathrm{O}$. м., когда система ее әлементарных делителей обладает следующими свойствами: 1) для $\lambda \neq \pm 1$ элементарные делители $(x-\lambda)^{m}$ и


\end{document}